\documentclass[12pt]{exam}
\usepackage{amsthm}
\usepackage{libertine}
\usepackage[utf8]{inputenc}
\usepackage[margin=1in]{geometry}
\usepackage{amsmath,amssymb}
\usepackage{multicol}
\usepackage[shortlabels]{enumitem}
\usepackage{siunitx}
\usepackage{cancel}
\usepackage{graphicx}
\usepackage{pgfplots}
\usepackage{listings}
\usepackage{tikz}


\pgfplotsset{width=10cm,compat=1.9}
\usepgfplotslibrary{external}
\tikzexternalize

\newcommand{\class}{APCS} % This is the name of the course 
\newcommand{\examnum}{HW1} % This is the name of the assignment
\newcommand{\examdate}{9-2-26} % This is the due date
\newcommand{\timelimit}{}





\begin{document}
\pagestyle{plain}
\thispagestyle{empty}

\noindent
\begin{tabular*}{\textwidth}{l @{\extracolsep{\fill}} r @{\extracolsep{6pt}} l}
\textbf{\class} & 
\textbf{Name:} & \textit{Toby Semple}\\ %Your name here instead, obviously 
\textbf{\examnum} &&\\
\textbf{\examdate} &&\\
\end{tabular*}\\
\rule[2ex]{\textwidth}{2pt}
% ---


\textbf{1:}
\begin{align*}
A &= \{1, 3, 12, 35\} \\
B &= \{3, 7, 12, 20\} \\
C &= \{x \mid x \text{ is a prime number}\}
\end{align*}

\begin{enumerate}
    \item[\textbf{a)}] $A \cap B = \{3, 12\}$
    \item[\textbf{b)}] $(A \cup B) \setminus C = \{1, 12, 20, 35\}$
    \item[\textbf{c)}] $A \cup (B \setminus C) = \{1, 3, 12, 20, 35\}$
\end{enumerate}

'a' and 'b' are both subsets of 'c'. There are no disjoints between these three sets.

\vspace{1em}

\textbf{8) a)}
\begin{align*}
(A \setminus B) \cap C &= (A \land \neg B) \land C \\
(A \cap C) \setminus B &= (A \land C) \land \neg B
\end{align*}
\begin{equation*}
(A \land \neg B) \land C = (A \land C) \land \neg B \qquad \text{commutative}
\end{equation*}

\textbf{b)}
\begin{align*}
(A \cap B) \setminus B &= (A \land B) \land \neg B \\
&= A \land (B \land \neg B) \qquad \text{commutative}
\end{align*}

The resulting contradiction shows the set is empty, as nothing can satisfy the entry requirements.

\textbf{c)}
\begin{align*}
A \setminus (A \setminus B) &= A \land \neg (A \land \neg B) \\
A \cap B &= A \land B
\end{align*}

\begin{align*}
& A \land \neg (A \land \neg B) & \\
&= A \land (\neg A \lor B) & \text{De Morgan} \\
&= (A \land \neg A) \lor (A \land B) & \text{Distributive} \\
&= \text{False} \lor (A \land B) & \text{Contradiction} \\
&= A \land B & \text{Identity}
\end{align*}

\begin{equation*}
A \land \neg (A \land \neg B) = A \land B
\end{equation*}

\newpage

\textbf{14}
\begin{align*}
B \Delta C \quad &\to \quad A \Delta (B \Delta C) \\
A \Delta B \quad &\to \quad (A \Delta B) \Delta C
\end{align*}

\vspace{1.5em}

\textbf{1.5: 2) a)}
\begin{align*}
P &= \text{Mary gets a good appartment Price} \\
A &= \text{Mary finds a nice appartment} \\
M &= \text{Mary moves out} \\
&(P \land A) \to M
\end{align*}

\textbf{b)}
\begin{align*}
C &= \text{good credit history} \\
D &= \text{adequate down payment} \\
M &= \text{you get a mortgage} \\
&M \to (C \land D)
\end{align*}

\textbf{c)}
\begin{align*}
S &= \text{Someone stops Jhon} \\
D &= \text{Jhon drops out} \\
&\neg S \to D
\end{align*}

\textbf{d)}
\begin{align*}
A &= x \text{ is dvsbl by } 4 \\
B &= x \text{ is dvsbl by } 6 \\
P &= x \text{ is prime} \\
&P \to \neg (A \lor B)
\end{align*}

\newpage

\textbf{8) a)}

The first phrase `$(P \to Q) \land (Q \to R)$' implies the transitive that can be found in the second, $(P \to R)$. Because the transitive $(P \to R)$ doesn't involve 'Q' at all, the two biconditionals in the second phrase are needed to cover the cases in which the conditionals are false.

\textbf{b)}

For the statement to not be a tautology:\\
\\
If 'P' is true, then 'Q' must be false to make the first conditional false. But if Q is false, then that means that $(Q \to R)$ is true by default.\\
\\
If 'P' is false, then the first conditional is true by default, making the second irrelevant.\\

\noindent\textbf{(a) $P \rightarrow (Q \rightarrow R)$}
\begin{table}[h]
    \centering 
    \begin{tabular}{ccccc}
        \textbf{P} & \textbf{Q} & \textbf{R} & \textbf{$Q \rightarrow R$} & \textbf{$P \rightarrow (Q \rightarrow R)$} \\
        T & T & T & T & T \\
        T & T & F & F & F \\
        T & F & T & T & T \\
        T & F & F & T & T \\
        F & T & T & T & T \\
        F & T & F & F & T \\
        F & F & T & T & T \\
        F & F & F & T & T \\
    \end{tabular}
\end{table}

\noindent\textbf{(b) $Q \rightarrow (P \rightarrow R)$}
\begin{table}[h]
    \centering 
    \begin{tabular}{ccccc}
        \textbf{P} & \textbf{Q} & \textbf{R} & \textbf{$P \rightarrow R$} & \textbf{$Q \rightarrow (P \rightarrow R)$} \\
        T & T & T & T & T \\
        T & T & F & F & F \\
        T & F & T & T & T \\
        T & F & F & F & T \\
        F & T & T & T & T \\
        F & T & F & T & T \\
        F & F & T & T & T \\
        F & F & F & T & T \\
    \end{tabular}
\end{table}

\noindent\textbf{(c) $(P \rightarrow Q) \land (P \rightarrow R)$}
\begin{table}[h]
    \centering 
    \begin{tabular}{cccccc}
        \textbf{P} & \textbf{Q} & \textbf{R} & \textbf{$P \rightarrow Q$} & \textbf{$P \rightarrow R$} & \textbf{$(P \rightarrow Q) \land (P \rightarrow R)$} \\
        T & T & T & T & T & T \\
        T & T & F & T & F & F \\
        T & F & T & F & T & F \\
        T & F & F & F & F & F \\
        F & T & T & T & T & T \\
        F & T & F & T & T & T \\
        F & F & T & T & T & T \\
        F & F & F & T & T & T \\
    \end{tabular}
\end{table}

\noindent\textbf{(d) $(P \land Q) \rightarrow R$}
\begin{table}[h]
    \centering 
    \begin{tabular}{ccccc}
        \textbf{P} & \textbf{Q} & \textbf{R} & \textbf{$P \land Q$} & \textbf{$(P \land Q) \rightarrow R$} \\
        T & T & T & T & T \\
        T & T & F & T & F \\
        T & F & T & F & T \\
        T & F & F & F & T \\
        F & T & T & F & T \\
        F & T & F & F & T \\
        F & F & T & F & T \\
        F & F & F & F & T \\
    \end{tabular}
\end{table}

\noindent\textbf{(e) $P \rightarrow (Q \land R)$}
\begin{table}[h]
    \centering 
    \begin{tabular}{ccccc}
        \textbf{P} & \textbf{Q} & \textbf{R} & \textbf{$Q \land R$} & \textbf{$P \rightarrow (Q \land R)$} \\
        T & T & T & T & T \\
        T & T & F & F & F \\
        T & F & T & F & F \\
        T & F & F & F & F \\
        F & T & T & T & T \\
        F & T & F & F & T \\
        F & F & T & F & T \\
        F & F & F & F & T \\
    \end{tabular}
\end{table}

Comparing the outputs of the truth tables yields two distinct sets of logically equivalent formulas:
\begin{itemize}
    \item  Statements \textbf{(a)}, \textbf{(b)}, and \textbf{(d)} are logically equivalent.
    \item  Statements \textbf{(c)} and \textbf{(e)} are logically equivalent.
\end{itemize}

\end{document}
▋